%% slides/es/main.tex -- charla general en espanol.  Primera version, 2026-08-11.
%%
%% NO es una traduccion de en/main.tex.  Es una charla distinta, para un publico
%% general que ha oido hablar de la IA y poco mas, y CON LIMITE DE TIEMPO:
%% 45 minutos como maximo.  Ver es/README.md.
%%
%% Registro: expositivo y positivo.  El mazo tecnico en ingles hace autocritica
%% deliberada (los dos fallos de correccion encontrados, el corpus de mutaciones);
%% AQUI NO.  Esto es un relato de por que el proyecto tiene sentido, no una
%% auditoria de sus defectos.  Honesto, pero no flagelante.
%%
%% Construir:  make es      (desde slides/)
%%
%% Las cifras vienen de ../shared/numbers.tex, igual que en el mazo ingles, para
%% que una cifra corregida no pueda pudrirse en uno de los dos.
%%
%% Sin marcadores \fid{}: `make parity' compara el esqueleto de frames de los dos
%% mazos y aqui no aplica -- esta charla no se corresponde frame a frame con
%% ninguna otra.

\documentclass[11pt,aspectratio=169]{beamer}
%% SIN babel.  Esta instalacion de TeX no tiene `texlive-lang-spanish': no hay
%% spanish.ldf y language.dat solo trae patrones ingleses, de modo que
%% \usepackage[spanish]{babel} aborta la compilacion.  El texto espanol se
%% compone igual de bien (utf8 + T1); lo unico que se pierde es la particion
%% de palabras a final de linea, que en ingles seria INCORRECTA.  Asi que se
%% desactiva la particion por completo: mejor una linea algo mas suelta que un
%% guion mal puesto.  Si se construye en una maquina con el paquete instalado,
%% reponer la linea de babel y quitar el bloque de penalizaciones.
\hyphenpenalty=10000
\exhyphenpenalty=10000
\tolerance=2000
\emergencystretch=3em
%% shared/preamble.tex -- packages, theme and macros common to EN and ES.
%%
%% LANGUAGE-NEUTRAL ONLY.  Nothing here may contain prose in either language:
%% the whole point of the shared/ directory is that a formula or a listing
%% corrected in the English deck cannot rot in the Spanish one.  Babel is
%% loaded by each main.tex, BEFORE this file, with its own language option.

\usetheme{Madrid}
\usecolortheme{seahorse}
\setbeamertemplate{navigation symbols}{}
\setbeamertemplate{caption}[numbered]
\setbeamerfont{frametitle}{size=\large}
\setbeamertemplate{itemize items}[circle]
\setbeamertemplate{footline}[frame number]

%% lmodern before fontenc: with T1 and plain Computer Modern, pdflatex falls
%% back to bitmap (Type 3) fonts and the deck looks blurred on a projector.
\usepackage{lmodern}
\usepackage[T1]{fontenc}
\usepackage[utf8]{inputenc}
\usepackage{amsmath,amssymb,amsthm}
\usepackage{fancyvrb}
\usepackage{upquote}
\usepackage{booktabs}
%% NOT enumitem.  It and beamer fight over \labelenumi and the loser is the
%% build ("TeX capacity exceeded, grouping levels=255" on the first enumerate).
%% The proof-reader emits `\begin{itemize}[leftmargin=...]', which needs it --
%% so `make proofs' strips that optional argument when it cuts the fragment.
%% Mechanical, and recorded here because the fragment then differs from the
%% document the reader wrote.
\usepackage{array}
\usepackage{graphicx}

%% amsbook's `proposition' is not among beamer's theorem environments, and the
%% proof-reader emits one.  Declare it here so shared/proofs/*.tex compiles.
\theoremstyle{plain}
\newtheorem{proposition}{Proposition}

%% The proof-reader's own helper: shrink an over-wide display to \linewidth.
\providecommand{\fit}[1]{%
  \resizebox{\ifdim\width>\linewidth\linewidth\else\width\fi}{!}{#1}}

%% --- the EN/ES parity marker -------------------------------------------
%% Every frame opens with \fid{<stable-id>}.  It typesets nothing; `make
%% parity' extracts the ids from both decks and diffs them, so a frame added,
%% dropped or reordered in one language is a build-time failure rather than a
%% discovery on stage.  See project_slides_beamer: the predictable failure of a
%% two-language deck is silent divergence.
\newcommand{\fid}[1]{}

%% --- code listings ------------------------------------------------------
%% Sources are plain .txt under shared/code/, extracted from the tree by
%% `make code' (shared/code/MANIFEST names file, source and line range), so a
%% listing on a slide is the source as it stands, not a transcription of it.
\DefineVerbatimEnvironment{vnbcode}{Verbatim}%
  {fontsize=\scriptsize,xleftmargin=1em,commandchars=\\\{\}}
\newcommand{\codefile}[2][\scriptsize]{%
  \VerbatimInput[fontsize=#1,xleftmargin=0.5em]{../shared/code/#2}}

%% --- small semantic macros ---------------------------------------------
\newcommand{\vnb}{\textsc{vnb}}
\newcommand{\imps}{\textsc{imps}}
\newcommand{\code}[1]{\texttt{\small #1}}
\newcommand{\qeds}[1]{\texttt{\small #1}}
\newcommand{\Set}{\mathbf{Set}}

%% shared/numbers.tex -- every quoted figure in either deck, in ONE place.
%%
%% A talk about a project in progress quotes numbers that move weekly.  They
%% are defined here, once, with the date and the command that produced them, so
%% that refreshing them before a talk is a single edit and so that the English
%% and Spanish decks cannot quote different figures.
%%
%% AS OF 2026-08-10.  Sources:
%%   \Nproven, \Nsupport, \Nstated, \Naxioms, \Ndefs, \Nassert, \Ncatalog
%%       -- reference/THEOREMS.md header, produced by (catalog)
%%   \Nmodzero, \Ntrustnone, \Nbills
%%       -- reference/PROOF-DEBT.md, produced by (proof-debt-ledger)
%%   \Nkernel   -- grep -c "define (pi-" primitive-inferences.scm
%%   \Nunexer   -- kernel-rules-audit, printed at the end of a library load
%%   \Nsuite    -- timeout 900 mit-scheme --quiet --load test-suite.scm
%%                 RUN IT ALONE.  Sharing the box with a ./prover process makes
%%                 it die partway through the library load having run no checks,
%%                 printing no SUMMARY and exiting 0 (seen 2026-08-10).  If the
%%                 SUMMARY line is absent, the number below is not a number.
%%   \Nmnp, \Nmnpctl, \Nmut -- test-suite-negative.scm
%%   \Nload     -- ./VNB cold boot, compiled tree

\newcommand{\Ncatalog}{1769}
\newcommand{\Nproven}{301}
\newcommand{\Nsupport}{610}
\newcommand{\Nstated}{860}
\newcommand{\Naxioms}{212}
\newcommand{\Ndefs}{510}
\newcommand{\Nassert}{138}

%% Counted at the end of 2026-08-10:
%%   grep -c "proven \*\*modulo 0\*\*" reference/PROOF-DEBT.md  -> 109
%%   grep -c "^### "                                    -> 301  (every proven theorem)
%%
%% \Ntrustnone is NOT a grep.  The legend at the head of PROOF-DEBT.md contains
%% the string "trust: none" itself, so grep -c reports one more than the truth.
%% The figure is the tally over *proof-debt* of bills whose
%% (debt-trust-level bill) is `none'.
%%
%% Its day, in nine steps: 54 -> 32 (five asserted structure laws PROVEN)
%% -> 32 again (three nary bridge converses stamped definitional -- correct,
%% and worth nothing alone, because a fourth axiom shadowed them in every bill)
%% -> 19 (that fourth) -> 19 again (integral-domain-cancel-zero PROVEN --
%% shadowed in turn by zz-is-integral-domain) -> 11 (zz-is-integral-domain)
%% -> 10 (integral-domain-nontrivial) -> 10 again (comm-monoid-is-monoid --
%% shadowed by its partner in the one bill they share) -> 9 (qq-is-ring)
%% -> 8 (nn-add-monoid-is-comm-monoid, that partner)
%% -> 4 (principal-ideal-membership stamped definitional).
%%
%% THREE of the ten steps moved nothing on their own.  The tier is the WORST
%% leaf, so discharging one pays only when it is the LAST unwarranted leaf of
%% the bills that name it -- and all three were nonetheless necessary, since
%% the leaf shadowing them had to go too.  Predict with the what-if
%% (scratchpad/nary-what-if.scm, scratchpad/cheap-whatif.scm) BEFORE spending
%% the effort: it was right every time today.
%%
%% The 4 that remain have NO shadowing left -- each is the sole unwarranted
%% leaf of its bill: zz-is-euclidean-ring (zz-bezout), rr-is-metric-space
%% (rr-complete), inf-subsets-is-set (two bills).
\newcommand{\Nmodzero}{109}     % unconditional bills -- the strongest report
\newcommand{\Ntrustnone}{4}     % bills whose weakest leaf has no warrant at all
\newcommand{\Nbills}{192}       % proofs carrying any debt at all (301 - 109)

\newcommand{\Nkernel}{50}
\newcommand{\Nunexer}{14}

\newcommand{\Nsuite}{791}       % passed; full run 2026-08-10, alone on the box
\newcommand{\Nsuitefail}{0}

\newcommand{\Nmnp}{12}          % false statements the corpus requires to stay open
\newcommand{\Nmnpctl}{7}        % paired controls that must close
\newcommand{\Nmut}{5}           % guards deleted on purpose, with predicted red entries

\newcommand{\Nload}{24}         % seconds, cold boot, compiled


\title[Modelos de lenguaje y matemáticas]%
      {Modelos de lenguaje y matemáticas}
\subtitle{Por qué hace falta un verificador, y en qué he empleado el tiempo}
\author{Javier Thayer}
\date{agosto de 2026}

\begin{document}

\begin{frame}
  \titlepage
\end{frame}

\begin{frame}{De qué voy a hablar}
  Tres preguntas, en este orden:
  \begin{enumerate}
    \item \textbf{¿Qué es} un modelo de lenguaje? Lo suficiente para entender
          en qué es bueno y en qué falla.
    \item \textbf{¿Por qué} han empezado a aparecer en la matemática?
    \item \textbf{¿Qué les falta}, y por qué lo que les falta tiene un nombre
          antiguo y preciso: una \emph{demostración formal}.
  \end{enumerate}
  \vfill
  Y al final, brevemente, en qué he empleado yo el tiempo.
  \vfill
  {\small No hace falta saber nada de informática. Ni de matemática, más allá
   de haber visto alguna demostración alguna vez.}
\end{frame}

\section{Qué es un modelo de lenguaje}

\begin{frame}{El problema que resuelve}
  Un modelo de lenguaje hace una sola cosa:
  \begin{quote}
    dado un texto, decir \textbf{cómo podría continuar}.
  \end{quote}
  \vfill
  \begin{center}
    \large
    ``El teorema de Pitágoras relaciona los lados de un \underline{\hspace{3em}}''
  \end{center}
  \vfill
  No responde ``triángulo'' porque lo tenga apuntado en algún sitio. Responde
  que, entre todas las continuaciones posibles, esa es con mucho la más
  probable \emph{en los textos que ha visto}.
  \vfill
  Todo lo demás --- contestar preguntas, escribir código, traducir --- es esa
  operación repetida.
\end{frame}

\begin{frame}{Primero: el texto se corta en piezas}
  El modelo no ve letras ni palabras, sino \textbf{tokens}: piezas de palabra,
  del tamaño que resulte más económico.
  \vfill
  \begin{center}
    \large
    \texttt{in\,|\,cre\,|\,íble\,|\,mente} \qquad
    \texttt{Pi\,|\,tá\,|\,goras} \qquad
    \texttt{el\,|\,teorema}
  \end{center}
  \vfill
  Las palabras frecuentes son un token; las raras se parten. Un texto español
  típico gasta algo más de un token por palabra.
  \vfill
  Esto no es un detalle de implementación inocente: el modelo razona sobre
  \emph{esas} piezas. Contar letras, por ejemplo, le resulta antinatural, porque
  las letras no son las unidades que ve.
\end{frame}

\begin{frame}{Segundo: predecir la pieza siguiente}
  Para cada posición, el modelo produce una \textbf{probabilidad para cada
  token posible} del vocabulario:
  \vfill
  \begin{center}
    \begin{tabular}{lr}
      \toprule
      \multicolumn{2}{c}{``\dots los lados de un ''}\\
      \midrule
      \texttt{triángulo}   & 0.94 \\
      \texttt{polígono}    & 0.02 \\
      \texttt{cuadrado}    & 0.01 \\
      \dots                & \dots \\
      \bottomrule
    \end{tabular}
  \end{center}
  \vfill
  Se elige uno, se añade al texto, y se vuelve a empezar. Un párrafo son
  cientos de repeticiones de este paso.
  \vfill
  {\small La elección no siempre es la más probable: se muestrea. Por eso el
   mismo prompt da respuestas distintas.}
\end{frame}

\begin{frame}{El entrenamiento, en una frase}
  El modelo tiene un número enorme de parámetros ajustables. Entrenarlo es:
  \begin{quote}
    recorrer una cantidad enorme de texto, y ajustar los parámetros para que
    las continuaciones que \emph{de hecho} ocurrieron resulten más probables.
  \end{quote}
  \vfill
  Nada más. No se le enseña gramática, ni aritmética, ni lógica; no hay
  reglas escritas por nadie. Todo lo que parece saber es un efecto lateral de
  hacer bien esa predicción.
  \vfill
  \begin{block}{Lo sorprendente, y conviene decirlo}
    Que de un objetivo tan pobre salga algo que sostiene una conversación
    técnica no estaba previsto por nadie. Es un hecho empírico, no una teoría.
  \end{block}
\end{frame}

\begin{frame}{Lo que un modelo de lenguaje NO es}
  \begin{itemize}
    \item \textbf{No es una base de datos.} No hay una tabla de hechos que
          consultar, ni un sitio donde esté escrito lo que responde.
    \item \textbf{No comprueba nada.} Nada en su funcionamiento distingue una
          continuación verdadera de una meramente verosímil.
    \item \textbf{No sabe lo que no sabe.} No hay una señal interna que diga
          ``de esto no tengo ni idea''.
  \end{itemize}
  \vfill
  De ahí el fallo característico, y no es un defecto que se vaya a arreglar
  puliendo: \textbf{produce texto bien formado, fluido y convincente, que puede
  ser falso}. Está optimizado para la fluidez. La corrección no aparece en
  ninguna parte del objetivo.
\end{frame}

\section{Por qué en matemáticas}

\begin{frame}{Lo que estas herramientas hacen bien}
  Y que resulta que el matemático usa todo el rato:
  \begin{itemize}
    \item \textbf{Lenguaje técnico.} Leen y escriben con soltura la prosa
          peculiar en que está escrita la matemática.
    \item \textbf{Analogía.} ``Esto se parece a aquello'' es su operación
          nativa, y es también como se buscan las demostraciones.
    \item \textbf{Recuperación.} Traen a la mano el lema que uno sabe que
          existe pero no recuerda dónde.
    \item \textbf{Código.} Escriben programas razonablemente bien, lo que
          convierte lo experimental en barato.
  \end{itemize}
  \vfill
  Ninguna de estas cosas es demostrar. Todas son parte del trabajo.
\end{frame}

\begin{frame}{Dónde se están usando ya}
  Sin entrar en detalles, el patrón se repite:
  \begin{itemize}
    \item redactar y reescribir: pasar de la idea al texto publicable;
    \item buscar en la literatura, que hace tiempo dejó de ser abarcable;
    \item sugerir conjeturas y contraejemplos a partir de casos calculados;
    \item automatizar los pasos rutinarios de una demostración larga ---
          los que todo el mundo dice ``análogamente'' para no escribir;
    \item ayudar a \emph{formalizar}: pasar una demostración de papel a un
          sistema que la comprueba. De esto hablo en un momento.
  \end{itemize}
  \vfill
  Lo que casi nunca se afirma, y conviene notarlo, es que hayan
  \emph{demostrado} algo por su cuenta.
\end{frame}

\begin{frame}{Por qué es atractivo}
  Buena parte del tiempo de quien hace matemáticas no se va en tener ideas.
  Se va en:
  \vfill
  \begin{center}
    comprobar los casos pequeños \quad$\cdot$\quad
    arrastrar hipótesis por doce páginas \quad$\cdot$\quad
    verificar que el lema citado dice de verdad lo que hace falta \quad$\cdot$\quad
    rehacer la cuenta con el signo corregido
  \end{center}
  \vfill
  Es trabajo mecánico que exige atención sostenida, y el error se cuela
  justamente ahí, no en la idea central.
  \vfill
  Un ayudante incansable para esa parte sería valioso aunque no tuviera una
  sola idea propia.
\end{frame}

\section{Lo plausible y lo correcto}

\begin{frame}{El problema, en una línea}
  \begin{center}
    \Large
    Una demostración que \emph{parece} correcta\\[0.5ex]
    y una demostración correcta\\[0.5ex]
    son objetos distintos.
  \end{center}
  \vfill
  Y un modelo de lenguaje está construido, literalmente, para producir el
  primero.
  \vfill
  Esto no es una objeción moral contra la herramienta. Es una descripción de
  qué es: la fluidez es el objetivo; la corrección no está en el objetivo.
\end{frame}

\begin{frame}{Un ejemplo que no necesita ordenadores}
  Considérese
  \begin{center}
    \Large $\; p(n) = n^{2} + n + 41 \;$
  \end{center}
  \vfill
  \begin{center}
    \begin{tabular}{lllll}
      \toprule
      $p(0)=41$ & $p(1)=43$ & $p(2)=47$ & $p(3)=53$ & $p(4)=61$ \\
      \multicolumn{5}{c}{\dots\ y así hasta $n=39$: \textbf{todos primos}} \\
      \bottomrule
    \end{tabular}
  \end{center}
  \vfill
  Cuarenta casos seguidos. Cualquiera apostaría.
  \vfill
  \begin{center}
    \large
    $p(40) = 40^{2} + 40 + 41 = 41 \cdot 41$
  \end{center}
  \vfill
  La evidencia era abrumadora y la conclusión, falsa. ``Se ve que sí'' no es un
  argumento --- y ``se ve que sí'' es exactamente lo que sabe hacer un modelo
  entrenado en parecerse a lo que ha visto.
\end{frame}

\begin{frame}{Por qué esto pesa más en matemáticas}
  En casi cualquier otro uso, un error se diluye: el resumen queda algo peor,
  la traducción algo torpe, y sigue sirviendo.
  \vfill
  En una demostración no. Un paso malo entre doscientos buenos no deja una
  demostración un poco peor: deja \textbf{ninguna demostración}.
  \vfill
  Y el paso malo es justamente el que nadie mira: el que parece rutina.
  \vfill
  \begin{block}{La pregunta operativa}
    No es ``¿me fío de esta herramienta?''. Es:\\
    \textbf{¿puedo comprobar lo que produce, sin fiarme de ella?}
  \end{block}
\end{frame}

\section{Rigor: los modelos formales}

\begin{frame}{Qué es una demostración formal}
  Una demostración ordinaria está escrita para convencer a una persona: dice
  ``análogamente'', ``es claro que'', ``sin pérdida de generalidad''.
  \vfill
  Una demostración \textbf{formal} no tiene ninguna de esas palabras. Es una
  sucesión de pasos en la que \emph{cada} paso está justificado por una regla
  fijada de antemano, de una lista corta.
  \vfill
  La consecuencia es la interesante:
  \begin{center}
    \large
    una máquina puede comprobarla, sin entender nada.
  \end{center}
  \vfill
  Comprobar es aplicar la lista de reglas y ver si encaja. Es aburrido,
  mecánico, y no admite opinión.
\end{frame}

\begin{frame}{La asimetría de la que vive todo esto}
  \begin{center}
    \Large
    Encontrar una demostración es difícil.\\[1ex]
    \emph{Comprobarla} es fácil.
  \end{center}
  \vfill
  Nadie sabe hacer lo primero de forma automática, en general. Lo segundo es
  un programa modesto, que además puede escribirse una vez y revisarse con
  cuidado.
  \vfill
  Esa asimetría es la que hace que la combinación tenga sentido. Si comprobar
  fuera tan difícil como descubrir, no habría nada que hacer aquí.
\end{frame}

\begin{frame}{La consecuencia: proponer y disponer}
  \begin{center}
    \large
    El modelo \textbf{propone}. \qquad El verificador \textbf{dispone}.
  \end{center}
  \vfill
  \begin{itemize}
    \item El modelo puede equivocarse cuanto quiera: sus propuestas no entran
          en la biblioteca hasta que pasan.
    \item Lo que hay que revisar deja de ser ``¿es de fiar el modelo?'' ---
          pregunta a la que nadie sabe responder --- y pasa a ser
          ``¿es correcto el verificador?'', que es un programa concreto y
          pequeño.
    \item Y el resultado es comprobable \emph{por terceros}, sin confiar ni en
          el modelo ni en quien lo usó.
  \end{itemize}
  \vfill
  La confianza se traslada de algo opaco y enorme a algo transparente y
  pequeño. Eso es todo lo que se pide.
\end{frame}

\begin{frame}{Qué se le pide entonces al verificador}
  Si va a ser la pieza en la que se confía, tiene que cumplir tres cosas:
  \vfill
  \begin{enumerate}
    \item \textbf{Núcleo pequeño.} Las reglas de inferencia de las que depende
          todo deben caber en una tarde de lectura.
    \item \textbf{Lenguaje escribible.} El modelo tiene que poder producir
          algo en ese lenguaje. Un formalismo que sólo un experto sabe teclear
          no sirve de nada aquí.
    \item \textbf{Cuentas claras.} Debe decir de qué depende cada teorema: qué
          se ha demostrado y qué se ha supuesto. Un sistema que no distingue
          las dos cosas no da rigor, lo simula.
  \end{enumerate}
  \vfill
  Las tres son exigencias de \emph{ingeniería}, no de lógica. Y son la razón
  de que me pusiera a construir uno.
\end{frame}

\section{VNB}

\begin{frame}{Qué he estado construyendo}
  \textbf{\vnb{}}: un verificador de demostraciones, y algo más que eso --- un
  asistente para \emph{encontrarlas}.
  \vfill
  \begin{itemize}
    \item Sobre teoría de conjuntos (von Neumann--Bernays), que es el lenguaje
          en el que ya está escrita la matemática ordinaria.
    \item Con las estructuras habituales --- un grupo, un anillo, un espacio
          métrico --- como lo que son: tuplas con propiedades.
    \item Y con una interfaz pensada para que escribir una demostración se
          parezca a hacer matemáticas, no a programar.
  \end{itemize}
  \vfill
  \begin{block}{El dato que más dice sobre el proyecto}
    Está escrito casi por completo \emph{por} un modelo de lenguaje, bajo mi
    dirección. Es a la vez la herramienta y el experimento.
  \end{block}
\end{frame}

\begin{frame}{Cada teorema viene con su factura}
  Cuando una demostración se cierra, el sistema no dice sólo ``correcto''.
  Dice \emph{de qué depende}:
  \vfill
  \begin{center}
    \code{demostrado módulo 0} \\[0.5ex]
    {\small nada supuesto: sólo las reglas del núcleo y las definiciones}
  \end{center}
  \vfill
  \begin{center}
    \code{demostrado módulo \{a, b\}} \\[0.5ex]
    {\small correcto \emph{si} se conceden $a$ y $b$, que nadie ha demostrado}
  \end{center}
  \vfill
  Parece un detalle contable y es lo contrario: es la diferencia entre un
  sistema que da rigor y uno que lo aparenta. De \Nproven{} teoremas de la
  biblioteca, \Nmodzero{} no suponen absolutamente nada.
\end{frame}

\begin{frame}{El estado, en tres cifras}
  \begin{center}
    \begin{tabular}{rl}
      \toprule
      \Nproven{} & teoremas demostrados y comprobados \\
      \Nmodzero{} & de ellos sin suponer nada en absoluto \\
      \Nload{} s & tarda la biblioteca entera en \emph{re}-demostrarse \\
      \bottomrule
    \end{tabular}
  \end{center}
  \vfill
  Esa última cifra es la que importa para el uso diario: nada se da por bueno
  porque ayer funcionara. Cada arranque vuelve a comprobarlo todo desde las
  reglas.
  \vfill
  {\small Análisis, álgebra, topología: continuidad y compacidad, el lema de
   Zorn, la irracionalidad de $\sqrt{2}$, el teorema del valor medio, la forma
   normal de Smith.}
\end{frame}

\begin{frame}{Lo que quiero averiguar}
  Lo interesante no es la biblioteca. Es esta pregunta:
  \vfill
  \begin{quote}
    \large
    ¿Basta con que uno escriba el \textbf{preámbulo} --- las definiciones, el
    enunciado, las ideas de la demostración en prosa --- para que la máquina
    cierre el resto?
  \end{quote}
  \vfill
  Sería la forma útil de todo esto: el matemático hace lo que sabe hacer
  --- tener la idea --- y no hace lo que hace mal --- la contabilidad de los
  detalles.
  \vfill
  \begin{block}{Qué contaría como respuesta}
    Un preámbulo para un teorema que yo \emph{no} haya demostrado a mano, y del
    que salga una demostración comprobada. Todavía no está.
  \end{block}
\end{frame}

\begin{frame}{Dónde creo que esto sirve antes}
  No en la investigación de frontera, sino donde el rigor es el contenido:
  \vfill
  \begin{itemize}
    \item \textbf{Enseñanza.} Un estudiante escribe un argumento y recibe, no
          una nota, sino el paso exacto que no se sostiene.
    \item \textbf{Cualquier razonamiento que convenga fijar} --- una
          especificación, un protocolo, un argumento que se va a repetir ---
          donde un modelo formal obliga a decir lo que se estaba suponiendo.
  \end{itemize}
  \vfill
  En ambos casos el valor no es que la máquina sea lista. Es que no deja pasar
  lo que no se ha dicho.
\end{frame}

\begin{frame}{En resumen}
  \begin{itemize}
    \item Un modelo de lenguaje predice la continuación de un texto. Todo lo
          que hace sale de ahí, incluidas sus virtudes y su fallo.
    \item Su fallo --- ser convincente sin ser correcto --- es tolerable en
          casi todo, y no lo es en una demostración.
    \item Existe desde hace mucho la noción que arregla eso: la demostración
          formal, comprobable por una máquina que no entiende nada.
    \item Y existe la asimetría que lo hace práctico: descubrir es difícil,
          comprobar es fácil.
  \end{itemize}
  \vfill
  \begin{center}
    \large
    En eso he empleado el tiempo: en construir la parte que comprueba,\\
    para poder usar sin miedo la parte que propone.
  \end{center}
\end{frame}

\begin{frame}
  \begin{center}
    \Huge Gracias
  \end{center}
  \vfill
  \begin{center}
    Preguntas
  \end{center}
\end{frame}

\end{document}
